\documentclass[12pt]{article}
\usepackage{geometry}
\geometry{a4paper}
\usepackage[round, sort]{natbib}
\usepackage{graphicx}
\usepackage[T1]{fontenc}
\usepackage[utf8]{inputenc}
\usepackage{textcomp}
\usepackage{gensymb}
\usepackage{amsmath}
\usepackage{amssymb}
\usepackage{authblk}
\usepackage[running]{lineno}
\usepackage{setspace}
\usepackage{longtable}
\usepackage{hyperref}
\usepackage{times}
%\usepackage{colorbox}


\topmargin 0.0cm
\oddsidemargin 0.2cm
\textwidth 16cm
\textheight 21cm
\footskip 1.0cm
\doublespacing

\renewcommand\Authfont{\fontsize{12}{14.4}\selectfont}
\renewcommand\Affilfont{\fontsize{10}{10.8}\itshape}




\title{\large{The impacts of geographic range size and niche breadth on species centrality in floral visitation networks}}


\author[a,*]{Grant Foster}
\author[a,b]{Cleber Ten Caten}
\author[a]{Tad A Dallas}


\affil[a]{Department of Biological Sciences, University of South Carolina, Columbia, SC, 29208 }
\affil[b]{ Department of Botany, University of Wisconsin, Madison, WI, 53706}   

\renewcommand\Authands{ and }
\date{ \small *Corresponding author: grant.t.foster500@gmail.com}



%% Potential reviewers
% Names and email addresses of potential reviewers
%
%
%
%
%
%



\begin{document}

%\maketitle

\setstretch{1}
\vspace{-1cm}

\noindent \large{\textbf{Title}: The impacts of geographic range size and niche breadth on species centrality in floral visitation networks} \\

\noindent \textbf{Running title}: Effects of Pollinator Range Size on Centrality \\

%\noindent \textbf{Author contributions}: \textbf{GF:} Conceptualization, Methodology, Formal analysis, Investigation, Writing - Original Draft, Visualization. \textbf{CTC:} Methodology, Formal analysis, Validation, Writing - Review \& Editing. \textbf{TAD:} Conceptualization, Investigation, Writing - Review \& Editing, Supervision.   \\

\noindent \textbf{Acknowledgments}: This work has been supported by ANONYMIZED \\%the University of South Carolina and the National Science Foundation (NSF-DEB-2420769 and NSF-DEB-2017826).  \\

%\noindent \textbf{Acknowledgments}: This work has been supported by the University of South Carolina and the National Science Foundation (NSF-DEB-2420769 and NSF-DEB-2017826).  \\

\noindent \textbf{Data accessibility}: Data and $R$ code necessary to recreate this analysis is available on on figshare at \hyperlink{https://figshare.com/s/ab0b46493fd69cf4a2c8}{https://figshare.com/s/ab0b46493fd69cf4a2c8}.\\

% TAD: don't give two different repos for the project which could conflict (if you randomly decided to change the github repo). 
% GF: Alrighty; I've removed the github one then and will just reference the figshare since it's got a doi


\noindent \textbf{Keywords}: Range Size, Ecological Networks, Interaction Turnover, Pollination, Floral-Visitor, Centrality, Macroecology \\

\noindent \textbf{Conflict of interest}: The authors have no conflicts of interest to declare.\\

\clearpage

\linenumbers


\noindent {\large The impacts of geographic range size and niche breadth on species centrality in floral visitation networks}





30330

\setstretch{1.3}
\subsection{Abstract}
Understanding the factors that influence species' roles within ecological networks can provide insights into the rules governing the assembly of interaction networks. Centrality, a property describing how connected a given node is to the rest of a network, is often used to identify species whose loss may result in disproportionate effects on biotic communities. Understanding the drivers of species' centrality in floral visitation networks is crucial for understanding network responses to global change. Here, we test how geographic range size and niche breadth influence species centrality in floral visitation networks across spatial scales. We combine a globally distributed dataset of floral visitation networks comprising 30,330 interactions among 6,857 species with large-scale occurrence data to evaluate how species’ geographic and environmental range sizes impact their role in networks across scales. In global metawebs we observe a strong positive effect of geographic range size on closeness centrality, but that effect largely disappeared after accounting for spatial turnover in interactions. Instead, niche breadth showed a positive relationship with the standardized centrality of flowering plant species in both network types. Site-level centrality was tightly linked to species ranges. Together, these results emphasize the need for a multiscale perspective to understand the drivers and consequences of changes in ecological network properties across environments and scales. 
%Older version-cut for PRSB word count. Understanding the factors that influence species' roles within ecological networks can provide insights into the rules governing the assembly of interaction networks. Centrality, a suite of measures describing how connected a given node is to the rest of a network, is often used to identify species whose loss may result in disproportionate effects on biotic communities. Understanding the drivers of species' centrality in floral visitation networks is crucial, as declines in pollinator abundance and diversity threaten to reduce pollination services globally. Here, we test how geographic range size and niche breadth influence species centrality in floral visitation networks across spatial scales. Using a globally distributed dataset of floral visitation networks comprising 31,677 interactions among 2,703 plant and 4,305 floral visitor species across 99 sites, we model species’ geographic and environmental ranges from large-scale occurrence data to evaluate how these traits shape species’ roles in both local networks and global metawebs. In global metawebs, geographic range size was positively related to closeness centrality. However, after accounting for spatial turnover in interactions, this effect largely disappeared or reversed. Instead, niche breadth showed a positive relationship with the standardized centrality of flowering plant species in both network types. Site-level centrality is less tightly linked to species ranges. For insect-pollinated plant species, local centrality was negatively related to geographic range size and positively related to both niche breadth and the interaction between the two, while both measures were unrelated to the local centrality of nodes within avian networks. Together, these results emphasize the need for a multiscale perspective to understand the drivers and consequences of changes in ecological network properties across environments and scales. 





%\subsection*{ESA Abstract}
%\paragraph{}
%In mutualistic interaction networks, a species’ role in a given network can be described using several mathematical properties. Node centrality-a suite of measures describing how connected a given node is to the rest of the network-is a commonly studied metric as removing highly central species from a network is likely to have outsized impacts on biotic communities. Understanding the drivers of species' centrality in floral visitation networks is crucial, as declines in pollinator abundance and diversity threaten to reduce pollination services globally. One potential driver of species centrality is the size of its range. We hypothesize that, at the level of the metaweb (the global set of observed interactions) species centrality in mutualistic networks would positively scale with both environmental and geographic range size, as wide-ranging species exhibit turnover in interaction partners across space or environments. We also expected this positive relationship to hold in the context of local interaction networks as well, as broad ranges may be associated with generalist partner preferences. Here we test these expectations by characterizing these associations in floral visitation networks. We used large-scale occurrence datasets to calculate the geographic and environmental ranges occupied by 2,703 plant and 4,305 floral visitor species, distributed across 99 globally distributed floral-visitation networks comprising 31,677 unique interactions.

%\paragraph{}
%We found that at the global metaweb level, a species’ closeness centrality in both insect-plant and avian-plant floral visitor metawebs was positively related to geographic range size but was negatively related to environmental breadth in only the insect metaweb. However, after using null model approaches to standardize for the effect of potential interaction turnover, standardized closeness showed no relationship with either measure of range size in the insect metaweb, and a negative relationship with environmental breadth in the avian network. At a local scale, centrality in insect-plant networks was positively related to environmental breadth but negatively related to geographic range size. In avian-plant networks, we observed a negative effect of environmental breadth on species centrality and no impact of geographic range size. Together, our results indicate that the impacts of range size on species centrality in large-scale metawebs are primarily driven by interaction turnover and the scaling of interaction specificity from local to global networks. Species with large geographic ranges may appear to be relatively generalist in global metawebs due to high interaction turnover, but local-scale interaction patterns may exhibit more nuanced relationships with species range limits.













\clearpage

\setstretch{1.3}
\subsection*{Introduction}


\paragraph*{}
%TS: Centrality is an important measure of node's contribution to a network, and understanding what makes species more central is important
Originating in mathematics, graph theory perspectives have become increasingly popular for the study of ecological networks throughout recent years \citep{bascompte2007plant, landi2018complexity, valdovinos2019mutualistic, guimaraes2020structure}. Under this framework, interacting species can be represented as a series of nodes connected by links (interactions). What constitutes an ``interaction" is variable, set by the context of the system, and may include predation events, infections \citep{dallas2017predicting}, floral visits \citep{Young21}, or any number of other biotic interactions. By quantifying the full set of links present in a given system, we have gained insights into the structural properties of species interaction networks that allow us to better predict their potential responses to change \citep{gilarranz2017effects}. For example, foundational work has linked network properties such as nestedness (the tendency for species with few interactions to link to species with many other interactions) and modularity (the tendency for species to interact more within given groups or ``modules'' than between them) to features such as network stability in response to perturbations or species extinctions \citep{okuyama2008network, thebault2010stability, suweis2013emergence}.




\paragraph*{}
%TS: Species Centrality in Empirical Networks is important, and people try to predict it
In addition to these network-level properties, we can also measure the contribution of an individual node to overall network structure considering node-level properties. One of the most commonly analyzed node-level properties is centrality: a suite of measures describing how connected a given node is to the rest of the network. Highly central species may represent keystone species in a network, and removal of these nodes is likely to have outsized impacts on their associated communities \citep{traveset2017plant, maia2019plant, martins2024propagation}. Identifying the eco-evolutionary processes that give rise to highly central species in local and regional-scale interaction networks is essential for understanding how ecological networks are organized and how they may respond to disturbances \citep{moulatlet2023species, martins2025determinants, zhang2025network}. Previous work has emphasized the importance of species traits \citep{watts2016influence, lazaro2020linking, maia2019plant} and evolutionary history \citep{Gonzales2015macroecology, burin2021macroevolutionary} on influencing both network and node-level properties. Investigating determinants of a species' role within floral-visitation networks (defined here as its position and the distribution of interactions relative to other members in a network, captured through node-level properties such as centrality), can provide insights into biological processes underpinning interaction network structure across scales. Understanding these drivers in the context of floral visitation networks, which are often characterized by interaction generalism \citep{waser1996generalization,stanton2003interacting} and frequent facultative interactions \citep{geib2012tracing}, is of particular conservation importance as global declines in pollinator abundance and diversity threaten to potentially reduce pollination services in many ecosystems \citep{potts2010global, rodger2021widespread}. 


%TS: Geographic range size may be associated with species role in a network. 
\paragraph*{}
In addition to morphological characteristics and phylogenetic history, the size of a species' geographic range may also impact its role in pollination networks at both local and global scales \citep{moulatlet2023species, dattilo2020species}. 
% TAD: what does "its" in the sentence above refer to? I see what you're doing with the transition, but it's pretty awk. 
% GF: Rewrote. Still has an "its", but hopefully more clear. 
As species ranges increase in size, so does the total number of potential interaction partners. As a consequence, in regional metawebs (the set of species interactions observed across multiple local webs \citep{poisot2012dissimilarity}) wide-ranging species may increase in their total number of interactions as both potential partners and interactions turn over across space \citep{fricke2020accelerating, martins2022global}. 
% TAD: so now we have 2 different flavors of metaweb, regional and global?
% GF: I'm hoping that I can use "regional" as a catch-all (to have a metaweb you have to choose an "region" [extent] to aggregate over. In our case that's the globe. 
This accumulation of interactions across area, parallel to the classical species-area relationship \citep{burkle2012shifts,sugiura2012scale, galiana2018spatial, dallas2021species}, may thus induce positive associations between species range size and species centrality at regional or global scales.
% TAD: but people have done interaction scaling, so why just be like "similar to species-area relationships"? You don't even have to cite my shit here if you don't want to, but I've done it. 
% GF: Hm; I was trying to reference the interaction-scaling lit here but maybe made things weird with the species-area relationship comparison. The citations here are mostly 
Additionally, species with broad geographic ranges may also be more central in local networks if geographic range size is associated with demographic or functional traits that may allow for generalist interaction strategies or high local abundance \citep{holt1997relationship, hurtado2024generalism, Moulatlet2025General}. This potential confounder makes it difficult to differentiate the potential reasons geographic range size may impact species roles in mutualistic networks. The development of null modeling approaches, in which many random permutations of network linkages are summarized into an observed distribution of potential network property values, has been instrumental in making meaningful comparisons across ecological networks of different size and connectance \citep{dormann2017identifying, guimaraes2020structure, bluthgen2024critical}. Here, we utilize a null modeling approach in order to tease apart the otherwise confounding effects of turnover and range size \textit{per se} on species centrality across scales. 



\paragraph*{}
% We can think of "range size" geographically or environmentally
Species' geographic range size is driven by a number of biogeographic factors, including evolutionary history \citep{webb2000geographic, alzate2025evolutionary}, dispersal limitation \citep{alzate2022understanding}, biotic interactions \citep{pigot2013species,galiana2023climate}, and climatic constraints \citep{hausharter2023niche}. As species increase in the geographic extent they occupy, they may geographically overlap with more potential partners, increasing their centrality. However, geographic range size might not be the strongest predictor of the number of potential partners. Instead, species that occur in more climatically diverse environments may co-occur with a larger set of potential partners \citep{slatyer2013niche, batstone2018using}. Therefore, species with broad abiotic niche breadths may have higher centrality in the global metaweb as a function of greater niche overlap with potential partners. A positive relationship between niche breadth and local centrality may also be observed if climatic and interaction generalism are positively associated, a feature observed in some food web systems \citep{slatyer2013niche}. Niche breadth and geographic range size are fundamentally linked, though the exact shape of the relationship is dependent on the geographic availability of environments across space \citep{dallas2025linking}.



%TS: Centrality means different things at the local and regional scale
\paragraph*{}
Biologically, species centrality has a fundamentally different meaning in local networks and global metawebs, and understanding how centrality scales between these contexts is still an open question in network ecology. Most existing work investigating the determinants of species centrality within networks to date has focused on local scales (\citep{maia2019plant}, but see \citep{moulatlet2023species}), as here node centrality has clear links to network properties \citep{gonzalez2010centrality, maia2019plant, cagua2019keystoneness}. Species that occur in many local networks and are highly central in many or all of them may also exhibit high centrality in regional or global metawebs (\citep{estrada2007characterization, martins2025determinants} though see \citep{dallas2022parasite}). However, even species with low local centrality can achieve high centrality in the global metaweb if they exhibit high interaction turnover across a wide geographic area \citep{carstensen2014beta, simanonok2014partitioning}. At a global level, these central species may appear as relative generalists with many potential partners. However, at local scales, the realized set of interactions of any given individual or population may be relatively narrow \citep{bison2015upscaling, arroyo2023intraspecific}. As such, understanding how centrality scales from local to global scales helps us better understand how population-level interaction patterns may vary across mutualistic species' ranges. 




% thesis
\paragraph*{}
To address these questions, we quantified variation in species-level centrality relative to geographic and environmental range size across a global assemblage of avian-plant and insect-plant floral visitation networks. Specifically, we tested whether species with broader geographic ranges or environmental niches occupy more central positions within local networks or global metawebs. We predicted that centrality will increase with range size and niche breadth in both contexts \ref{fig:concept}. At the metaweb scale, we expected species with broad ranges to turnover in interactions across geographic or environmental space, allowing them to accumulate more interactions across sites. At the local scale, we predicted a similar positive relationship driven by the positive association between interaction generalism and climatic or habitat generalism. Finally, we investigated whether these relationships persist after using null-model approaches to account for the impact of species co-occurrence on potential centrality measures. If range size or niche breadth continues to be positively associated with centrality under null models that control for co-occurrence, it would suggest a fundamental link between geographic or climate distributions and factors determining species role within mutualistic networks. Together, our results suggest that the impacts of range size on species centrality in large-scale metawebs are primarily driven by interaction turnover and the scaling of interaction specificity from local to global networks. When variation in interaction turnover is incorporated, niche breadth exerts the strongest positive effect on plant centrality within global networks. Consequently, species with large geographic ranges may appear to be relatively generalist in global metawebs due to high interaction turnover, but local-scale interaction patterns may exhibit more nuanced and variable relationships with species range limits. 


% \vspace{30pt}
% RQ1: Do we see a positive relationship between range size and species centrality in local or regional metawebs?
%     H1: As geographic range size increases, metaweb centrality will increase due, at least partially due to an increase in total number of connections.
%     H2: As geographic range size increases, local centrality will also increase due to features such as generalist species being associated with larger ranges
% RQ2: If so, does this appear when we use null models that account for potential co-occurrence?
%     H1: Even if we account for the effect of increased coocurence that can drive centrality in the metaweb, do we still see effects of geographic range size. 

%----------
\subsection*{Methods}

\paragraph*{Data}
We queried floral visitation networks through the Web of Life dataset, resulting in 97 networks with a total of 30,330 unique interactions between 6,857 species. We combined these networks into two global metawebs: one describing avian-plant interactions and another describing insect-plant interactions. Most studies only characterized one of these types of interactions; for the three studies that reported both insect and avian floral visits, we divided the network as appropriate. The avian-plant metaweb is comprised of 2,033 unique species-level interactions between 605 plant and 120 bird species across 31 sites. The insect-plant metaweb is larger, consisting of 28,297 interactions between 1,960 plants and 4,182 insect species across 66 sites. While expert range maps are available for most birds and some well-studied plant taxa, the majority of species in our study do not have up-to-date expert species range estimates available. In order to investigate the role of range size in determining species-level network traits, we estimated species ranges from available occurrence points. 

\paragraph*{Estimating Geographic Range Size}
For each species, we queried global occurrence records from the Global Biotic Information Facility (GBIF) using the \texttt{gbifdb} $R$ package. Occurrence points were first processed using the \texttt{coordinateCleaner} package; points were removed if they were present only in the fossil record, occurred within 2 km of country centroids, were not terrestrial, or were marked as an ``introduced" individual. To avoid overestimating species range size through including spurious or transient records in GBIF, we used a quantiled minimum convex polygon (MCP) approach. Within each continent -- based on a 7 continent model (Asia, Europe, Africa, Oceania, North America, Antarctica, and South America) -- we created the smallest minimum convex polygon that contains the central 95\% of the occurrence points within each continent. Range polygons were not created for species-continent combinations with less than 10 occurrence points. Total species' geographic range size was calculated as the total terrestrial area occupied across all continents. The points included within these polygons were also used for the calculation of environmental range (below). 2,554 species within our network had insufficient occurrence data to create species ranges; leaving a total of 4,303 species within our analysis (110 birds, 2,126 insects, and 2,067 plants). 

\paragraph{Estimating Niche Breadth} 
For each occurrence point, we extracted 56 climatic values from WorldClim at a 2.5 arc-minute resolution \citep{fick2017worldclim}. These values represent a number of climatic variables potentially important for limiting species' distributions, averaged across the years 1970-2000. As these variables represent a high-dimensional environmental space that is collinear in many aspects, we performed a principal component analysis in order to reduce the dimensionality of our niche space \citep{kriticos2014extending, dallas2025linking}. The first two principal components of this climatic space explain over 77\% of the total climatic variation. For each species, we plotted each occurrence point that fell within the 95\% quantiled geographic range polygon calculated above onto the Cartesian plane defined by these two PCA axes, combining points from all continents on which a species occurred. We then again employed a minimum convex polygon approach to draw an environmental range in this PCA space for each species (i.e., its niche breadth).

% TAD: you didn't do the 95% of points with extreme values removed or whatever for this one? 
% GF: We used to use all the points directly and repeat the 95% quantiled approach, but after talking about it we now just use the same 95% points that are used to calculate the geographic range. I changed the above text to try and help specify this (...we plotted each occurence point that fell within the 95\% quantiles geographic range polygon calculated above...)


\paragraph*{Centrality Measures}
Insect and avian floral visitor metawebs contained 6142 and 725 total species, respectively. The vast majority of species in each network were connected to one major component (the largest set of interconnected nodes). As centrality measures become non-comparable in unconnected graphs, we dropped species that were not part of these major components. This resulted in the removal of 51 (0.83\%) and 4 (0.55\%) species from the insect and avian metawebs, respectively. Closeness centrality (which describes the mean proximity of one node to all other reachable nodes within the network, Figure \ref{fig:concept}) was calculated for nodes in the remaining simply connected networks using the \texttt{igraph} package. 
% TAD: why give a weirdly technical def of closeness? It's not wrong or bad, but you want plant-frugivore folks reading this. It's proportional to the mean proximity of one node to all other reachable nodes in the network. GF: Casual-ified! 
We chose closeness centrality as an intermediate between hyper-local measures of centrality (e.g., degree centrality) and hyper-global measures (e.g., eigenvalue centrality). At the global metaweb scale, the rank of alternative measures of centrality were positively correlated, though to varying degrees (Figure \ref{fig:S2}).

% TAD: replace all 'metrics' with 'measures'. Metrics implies distance
%GF: I'm not certain this is true, but I replaced them anyway lol. But what about business "performance metrics"? 

\paragraph*{Null model expectations of species centrality} 
While centrality measures are designed to be comparable between nodes of the same network, their absolute value is heavily influenced by network size and connectance, potentially hindering comparisons across networks. To address this when measuring species-level centrality in local networks, we used a null-model approach to standardize centrality within each local network. For each local network, we first measured the closeness centrality of all component species species using a similar approach to the global metaweb.
% TAD: you use lots of possessives (e.g., 'its') and it makes it read weird. If that's just your style, stick to it, but I am not a fan. 
%GF: Took out a couple of the most confusing ones; hopefully the remaining ones are relatively clear. 
We then used the \texttt{vegan} package to generate 1000 randomized versions of each local network by permuting interaction links while preserving each node's degree (effectively preserving marginal totals of the interaction matrix). This allowed us to characterize the potential distribution of centrality values for each node, given that node's degree and the network's overall structure. We then $z$-transformed the observed centrality values using the mean and standard deviation of this null distribution to create a $z-$transformed measure of local closeness centrality for each species-site combination. 


\paragraph{} To extend this approach to the metaweb level, we pooled permuted local webs to construct null global metawebs. This allowed us to generate a null distribution of species-level closeness centrality across the entire metaweb, accounting for each species' presence across multiple local networks. Observed metaweb centrality values were then $z-$transformed relative to this distribution. We investigated both measured and $z-$transformed centrality at the metaweb scale. At the local scale we only modeled $z-$transformed centrality, as raw measures of centrality are not directly comparable across differently sized networks. One insect-plant and three avian-plant networks were too small to create realistic null distributions, and were subsequently excluded for local centrality analyses. 


\paragraph*{Statistical Analyses}
We examined the role of species range size and niche breadth on species centrality through Bayesian generalized linear models implemented in Stan\citep{carpenter2017stan} through the $R$ package \texttt{brms} \citep{burkner2017brms}. In order to facilitate comparison and meaningful interpretation of interaction effects, both geographic range sizes and niche breadth values were log-transformed and then rescaled to a mean of 0 and standard deviation of 1. All models therefore shared the generalized form of predicting a given species' centrality as a function of geographic range size, niche breadth, and the interaction between these predictors. Niche breadth or geographic range size may exert fundamentally different effects on the centrality of flowering plant species compared to floral visitors. As such, we analyze node types (floral visitor or flowering plant) separately for each network or set of networks, effectively repeating our modeling approach for all factorial combinations of: scale (metaweb vs local webs), network type (insect-plant vs avian-plant), and interactor type (floral visitor vs flowering plant). Accounting for our two alternative centrality measures at the metaweb scale (measured and $z-$transformed closeness centrality), this resulted in 12 models total. We use vague priors centered on 0 for all slope and intercept coefficients. All models were fit using 5 Markov chain Monte-Carlo simulations, each with 5,000 iterations following a 1,000 iteration burn-in period; posterior samples were sampled without thinning \citep{link2012thinning}. Simulation convergence was checked through visual inspection of traceplots. Gelman-Rubin statistics fell below 1.01 for all parameters across all models.  For slope terms we sampled the resultant posterior distributions and approximated the two-sided posterior probability ($POP_2$) as twice the proportion of posterior samples falling above or below 0 (whichever is smallest), a Bayesian analog to traditional frequentist p-values \citep{shi2021reconnecting}. Relationships between fit slope values across models are presented as pairs plots in supplemental figures \ref{fig:S6}-\ref{fig:S11}. 

\paragraph{}
For the subset of species which occurred across multiple sites, we further investigated the role interaction turnover may play in determining species centrality in global metawebs. As both the degree of turnover and number of sites at which a species occurs may have nonlinear effects on metaweb centrality, we investigate these factors through a generalized additive modeling (GAM) approach. For each node the that occurred in multiple sites (1,166 species), we calculated the Simpson's dissimilarity of interaction partner sets across all sites in which they occurred \citep{baselga2015comparing}, as implemented in the $R$ package \texttt{betapart} \citep{baselga2012betapart}. We created GAM models predicting metaweb centrality (either raw values or $z-$transformed, scaled to a mean of 0 and standard deviation of 1 to facilitate comparisons across models) as a function of the tensor product of interaction turnover and number of local sites. We investigate the $p-$value of tensor terms, as well as compare estimated degrees of freedom to maximum allowable degrees of freedom ($k'$) as well as visually inspected model residuals to assess model fit. %Plots illustrating the resultant partial effects of turnover across are show in Figure \ref{fig:gam}. 

% TAD: have you defined k'? might be worth just spelling it out instead of the notation. You want the reader to follow, and you've already been talking about tensor products. 
% GF: Definition added


% -------------------
\subsection*{Results}

\paragraph*{Species Occurrences and Interactions}
Species' geographic range sizes in insect-plant networks varied widely, ranging from 2,232 km$^2$ (\textit{Chaetanthera lycopodioides}, a range-restricted montane Aster endemic to Chile and Argentina \cite{arroyo2007display}) to over 119 million km$^2$ (\textit{Sonchus oleraceus}, a cosmopolitan weed native to Eurasia), with a median value of 3,749,391 km$^2$. Geographic distributions of species in avian-plant visitation networks ranged from 2,518 km$^2$ (\textit{Barbacenia longiflora}, a monocot forb restricted to Minas Gerais, Brazil) to over 85 million km$^2$ (\textit{Prunella vulgaris}, a cosmopolitan invasive weed), with a median value of 1,536,647 km$^2$. Across sites, geographic and environmental range size were positively correlated ($r$=0.667, $p <$ 0.001). Typical of large-scale interaction metawebs, realized connectance (the percentage of realized links out of the total possible number of interactions) was low \citep{jordano1987patterns}. At the global scale, only 0.32\% and 2.8\% of potential links were realized in insect and avian interaction metawebs, respectively. At the site-level, average connectance of insect-plant networks was 0.194, whereas average avian-plant network was 0.356. 

% TAD: numbers don't go in math font unless you're willing to put every number in math font. 
%GF: Fixed

%TD: what the fuck is POP_2? 
%GF: I put the explanation of this at the very end of the methods, but it is a weird metric. It just seemed like a better for shorthand for "our 95% Credible Interval does or doesn't overlap 0" as well as being a closer analog to like hypothesis testing. But keeping it isn't a make or break for me. 
% TD: eh. Give it a cooler name? \omega_{95}? Idk man. 


%Effect in metaweb
\paragraph*{Global Metawebs}
At the metaweb scale, geographic range size was positively associated with species centrality of all node types across both insect and avian networks (Figure \ref{fig:MetawebEffect}). However, once accounting for potential co-occurrence through $z$-transformation, geographic range size showed no association with metaweb centrality for avian or insect floral visitors, or insect-visited plant species, and instead showed a negative association with the centrality of avian-visited plant species ($POP_2=$0.020). Niche breadth had a negative impact on centrality for plants and insects within the insect-plant metaweb after controlling for geographic range size ($POP_2<$0.001, $POP_2=$0.028). After $z$-transformation, niche breadth instead had a positive impact only on plant centrality in both insect and avian metawebs ($POP_2<$0.001 across both networks) while showing no effect for floral visitors. There was no significant interaction effect between geographic range size and niche breadth on unstandardized metaweb centrality across node types. However, after $z$-transformation, plants in both network types showed positive interaction effects ($POP_2<$0.001 and $POP_2=0.009$ for insect and bird-visited plants, respectively).



\paragraph{} For species that occurred across multiple sites, interaction turnover varied from 0 (complete conservation or nestedness of interactions across sites) to 1 (no overlap in interactions across sites), with a mean value of 0.644. The tensor product of interaction turnover and number of sites was statistically significant in predicting both raw and $z-$transformed metaweb centrality ($p<0.001$), but varied widely in the amount of variance explained. For the GAM model predicting measured metaweb centrality $R^2=0.231$, with interaction turnover exhibiting a positive partial effect across much of the realized parameter space, particularly for species with the highest centrality values (Figure \ref{fig:gam}). After $z$-transformation, the tensor product of turnover and number of sites still exhibited a statistically significant effect ($p<0.001$) but explained little of the overall variation ($R^2=0.054$). Taken together these results further suggest that positive relationships between species' geographic range size and measured metaweb centrality are largely driven by interaction turnover across sites. 



%Effect on local webs
\paragraph*{Local Webs}
Average species centrality in local floral visitation networks was positively correlated with their centrality in their respective global metawebs ($r=0.480$ and $p<0.001$ for insect webs, while $r=0.079$ and $p=0.035$ for avian webs). Species that occurred at multiple sites showed variable centrality across their respective local webs (Figure \ref{fig:S3}). Geographic range size was negatively related to the local centrality of insect-visited plant species (Figure \ref{fig:LocalPlot}, $POP_2=$0.018). Conversely, niche breadth was positively related to the average local centrality of nodes within insect networks ($POP_2<$0.001 for plants, $POP_2=$0.004 for insects). Insect-visited plants also exhibited a positive interaction term between geographic range size and niche breadth on average local centrality ($POP_2=$0.022). We failed to find any significant effect of range size, niche breadth, or their interaction on the local centrality of species within avian networks.







\subsection*{Discussion}
%The effect of geographic range size on metaweb centrality is mediated through cooccurrence. 
\paragraph*{} Generally we find species' geographic range size and niche breadth are more clearly linked to their centrality in aggregate plant-pollinator metawebs than within local interaction networks. The positive impact of geographic range size on metaweb centrality we observed across all node classes may be caused by a number of potential drivers. Species with broad geographic ranges are likely to co-occur with more potential interaction partners across their range (Figure \ref{fig:S4}). In this case, even if species are relatively peripheral across local networks, they may be relatively central at the metaweb level if interaction turnover is high across their range \citep{martins2025determinants}. Our results support this interpretation, as while species with large geographic ranges tended to be more central in the global metaweb, they did not tend to be more central than null expectations based on their connections across all local webs (Figure \ref{fig:MetawebEffect}). Additionally, generalized additive models indicated a large partial effect of interaction turnover on measured metaweb centrality, but this affect largely disappeared once these centrality values were $z-$transformed (Figure \ref{fig:gam}). Together, these patterns suggest that, at the global scale, species centrality in these metawebs is shaped more by neutral processes-namely the accumulation of interactions across space-than by the specific identity of those interactions. 

\paragraph*{}
Contrary to our expectations, standardized centrality was negatively related to geographic range size of plants in avian-plant networks after accounting for niche breadth. This pattern may indicate slower accumulation of interactions across space (lower interaction turnover) than otherwise expected. This pattern may be driven by wide-ranging, disturbance-adapted species that occur in many networks but are less likely to form interactions with narrowly distributed specialists. Nectarivory has evolved numerous times within birds, but contemporary biogeographical distributions of avian nectarivore diversity is clustered in the neotropics and parts of Australasia \citep{kissling2012bird}. This spatial clustering of potential partner diversity may drive lower rates of interaction turnover of avian-pollinated plants across large spatial scales. 

%The effect of environmental range size only impacts plants (mostly when $z-$transformed). 
\paragraph*{}
After $z$-transformation, metaweb centrality was positively related to niche breadth in plants of both metawebs, but was unrelated in floral visitors. If plant species are limited in their dispersal by the availability of mutualists \citep{warren2014mutualism, duffy2017specialized}, species with broad ranges of potential partners may be able to exploit a more climatically diverse set of environments. The availability or efficacy of mutualists may increase the abiotic tolerances of some species \citep{warren2014mutualism, fowler2023geographic}, effectively expanding their climatic niche breadth. Batstone et al. (2018) outline how environmental variability may favor interaction generalism by altering the availability or efficacy of mutualists \citep{batstone2018using}. Similar mechanisms may occur across an organism's range; plants with climatically diverse ranges may experience parallel changes in partner availability or efficacy, promoting interaction generalism \citep{hurtado2024generalism}. Alternatively, in systems characterized by widespread generalism, plant species with broad niche breadths may participate in more interactions as they co-occur with a broader complement of potential pollinators. Understanding the mechanisms of this interplay between a species' climatic niche breadth and its position within local and global networks warrants further study. 

\paragraph*{}
In contrast to global metawebs, both geographic range and niche breadth show a reduced effect on species centrality at the local scale. In insect networks, average local centrality was positively related to niche breadth across both node types, but negatively related to geographic range size in insect-visited plant species. We also detected a positive interaction between range size and niche breadth on the centrality of plant species in insect-floral visitation networks, indicating that their combined influence shapes local structure of these networks (Figure \ref{fig:LocalPlot}). Species with broad geographic ranges tended to be more central when they also had broader climatic niches, whereas species with narrow ranges were more central when they exhibited broader niche breadths. This suggests multiple ecological pathways or strategies that may result in high centrality in local networks. Widely distributed species with broad abiotic tolerances may be relatively generalized in their interaction patterns at local scales, increasing centrality \citep{hurtado2024generalism}. Alternatively, narrowly distributed climatic specialists may also be highly central, potentially through achieving high abundance in restricted environments \citep{suarez2023dispersal}. The functional consequences of a species' centrality within a network may vary greatly across local, regional, and global scales, and further investigation of how centrality scales across levels of organization is of great interest. In our study, networks are either analyzed at the site level, or aggregated into global metawebs. Fully understanding the processes of scaling from local to global requires hierarchical sampling designs replicated across large geographic extents. While understanding the scaling of network properties across space, environments, and time has received a great amount of conceptual attention in recent years \citep{tylianakis2008global,morales2008effect, burkle2011future,  moulatlet2023species}, the logistical challenges of research across these scales has limited our ability to fully explore these questions in some cases. 

\paragraph*{}
The emergent consequences of a species' centrality may ultimately differ across local to global scales. For example, at local scales the perturbation of highly central species may propagate farther or more strongly through a mutualistic community than the same perturbation applied to more peripheral species \citep{Cauga2019, delmas2019analysing, martins2024propagation}. Across larger scales, highly central species may play key roles in allowing mutualistic metacommunities to persist \citep{filotas2010effect}, especially in the face of habitat heterogeneity or disturbance \citep{prakash2004habitat}. Patterns of interaction turnover across space are closely linked to both the causes and consequences of centrality scaling relationships. While the temporal turnover of mutualistic interactions is fairly well appreciated in phenological contexts \citep{olesen2008temporal, chavez2020drivers}, empirical patterns of interaction turnover across space, habitats, and environments warrant further study \citep{tylianakis2017ecological}. 


\paragraph*{}
Understanding the factors that determine species positions in mutualistic networks across scales is a fundamental question in network ecology. Our study identifies geographic range size as a key determinant of species centrality in floral visitation metawebs, largely driven by the spatial accumulation of interactions across sites. At the local scale, we also find that climatic niche breadth interacts with geographic range size to predict plant centrality in insect-plant networks, highlighting the multiple pathways through species may become highly central at the site level. Together, these results emphasize the need for a multiscale perspective to understand the drivers and consequences of changes in ecological network properties across environments. 




\clearpage
\bibliographystyle{unsrt}
\bibliography{annot}

\clearpage
\subsection*{Figures}



\begin{figure}[h!]
    \includegraphics[width=24.7cm, angle=90]{figs/pdf/Concept_v2_busy_sized.pdf}
    \label{fig:concept}
\end{figure}

\addtocounter{figure}{0}
\begin{figure} [h!]
  \caption{Local interaction networks are embedded along natural environmental gradients (A), with species occurring in one or more sites and contributing interactions to a global metaweb, defined as the union of all observed interactions across sites. At the metaweb scale, species with broad geographic ranges or wide abiotic tolerances may accumulate interactions across multiple sites. At the local scale, such species may also be more central if interaction generalism covaries with environmental or habitat generalism. Species centrality can be quantified within individual site-level networks as well as in the global metaweb (node color). However, metaweb centrality is jointly influenced by the identity of interactions and by interaction turnover across sites. To isolate these effects, we implement a site-level null randomization (B) that shuffles interaction identities within each site while preserving the total number of interactions per species. Repeating this procedure generates null distributions of expected centrality values for each species at both local and metaweb scales (C), given the constraints of observed local network structure. Empirical centrality values can then be standardized (e.g., $z-$transformed) relative to these expectations, allowing scale-dependent effects of range size on network centrality to be evaluated (D).}
\end{figure}

\clearpage 




\begin{figure}[h!]
\centering
    \includegraphics[width=\columnwidth]{figs/pdf/Effectplot_Metaweb_Updates.pdf}
    \caption{Sampled posterior slope densities for multivariate Bayesian generalized linear models of species centrality in the global metaweb. Colored areas represent 95\% credible intervals of slope terms of the effect of niche breadth (top row, thermometer icon), geographic range size (center row, globe icon), and their interaction term (bottom row, chain icon).  Numbers and asterisks indicate approximate two-sided posterior probability $<0.05.$ Left columns represent models of measured closeness centrality as the response variable; right column represents models predicting centrality values $z-$transformed based on null randomization. After accounting for niche breadth, geographic range size has a strong positive effect of species centrality, but this effect disappears when z-standardized. In plant species, both niche breadth and the interaction between niche breadth and geographic range size positively impact $z$-transformed metaweb centrality.
    }
    \label{fig:MetawebEffect}
\end{figure}

\clearpage 

\begin{figure}[h!]
\centering
    \includegraphics[width=\columnwidth]{figs/pdf/gamplot.pdf}%
    \caption{Partial effects of interaction turnover on measured metaweb centrality (left) or $z-$transformed metaweb centrality (right) response values and number of unique local sites, according to fit tensor product GAMs. For comparison, both responses are rescaled to mean of 0 and standard deviation of one. Across the observed range of centrality values, the tensor product of interaction turnover and number of local sites occupied by a species had a positive effect of global metaweb centrality ($p<0.001, R^2=0.231$), especially for the most highly central species. In contrast, after $z-$transformation while the tensor product of turnover and number of local sites a statistically significant negative effect, this relationship explained less than 6\% of variation in metaweb centrality ($p<0.001, R^2=0.054$).}
    \label{fig:gam}
\end{figure}
\clearpage 

\begin{figure}[h!]
\centering
    \includegraphics[width=6in]{figs/pdf/LocalPlot_Composite_v2.pdf}%
    \caption{Relationships between average z-standardized centrality in local networks, geographic range size, and niche breadth. Species-level closeness centrality ($z-$transformed) in local interactions networks as a function of geographic range size (A,B) or environmental niche breadth (C,D). Panels (E-F) show sampled posterior slope densities for multivariate Bayesian generalized linear models of average species $z-$transformed closeness centrality in local metawebs. Colored areas represent 95\% credible intervals; numbers and asterisks indicate approximate two-sided posterior probability $<0.05.$ Niche breadth had a positive effect on the local centrality only of both node types within insect networks (E). Additionally, we also observe a positive interaction between niche breadth and geographic range size on standardized centrality for these node types as well(G).}
        \label{fig:LocalPlot}
\end{figure}

\clearpage 

\clearpage

\newcommand{\beginsupplement}{%
        \setcounter{page}{1}
        \setcounter{table}{0}
        \renewcommand{\thetable}{S\arabic{table}}%
        \setcounter{figure}{0}
        \renewcommand{\thefigure}{S\arabic{figure}}%
        }

\section*{Supplementary materials}
\beginsupplement

Floral visitation networks used in the dataset, queried from the Web of Life dataset, spanned different geographic extents depending on taxonomic focus. Our Plant-insect interaction networks (N=66) spanned the globe, ranging in size from 3 to over 15,000 interactions \cite{robertson1929flowers}. Avian-plant interaction networks (N=31) were exclusively sampled from the neotropics, ranging from solitary pairs of interactions to networks with as many as 248 unique interactions. Networks with less than four interactions were included in the global metaweb, but were excluded from analyses of local network properties, resulting in the exclusion of one insect and three avian networks from the analysis of local network properties. 

\begin{figure}[h!]
  \begin{center}
    \includegraphics[width=\textwidth]{figs/pdf/sup/sitemap.pdf}
    \caption{Map of networks used in our study. Each point represents a unique local network, created by observing avian (Red, N=31), or insect (Blue, N=66) floral visitations.}
    \label{fig:S1}
  \end{center}
\end{figure}
\newpage

\begin{figure}[h!]
  \begin{center}
    \includegraphics[width=\textwidth]{figs/pdf/sup/CentralityCorrelations.pdf}
    \caption{Pairwise correlation of common centrality measure as applied to our global metawebs. While the pearson's correlation of most measures are only slightly positive, spearman's rank correlation shows the ordering of nodes according to alternative measure is relatively conserved. In maintext we present results of closeness centrality, a middle-ground between hyper-local and hyper-global measures.}
    \label{fig:S2}
  \end{center}
\end{figure}

Species that occur in multiple local webs may exhibit different centrality values, dependent both on the compositions and connections observed within any given local community. For species that occur in multiple webs of our dataset, we attempt to predict mean $z-$transformed centrality across all webs where a species is found. To show the variability species may exhibit across networks, we plot the distributions of closeness centrality for species that occur at the most sites across our dataset. Differences between raw and z-standardized metaweb centrality stem from the number of couccuring extraguild partner species. Narrow-ranged plants that cooccur a small set of potential pollinator species are less likely to be highly central in the global metaweb, regardless of local connectivity. We show the relationship between geographic range size and the number of extraguild cooccuring species in \ref{fig:S4}. 


\begin{figure}[h!]
  \begin{center}
    \includegraphics[width=\textwidth]{figs/pdf/sup/Local_Centrality_Distributions.pdf}
    \caption{Distribution of closeness centrality in local networks for species that occur at multiple sites. Top row of plots represent nodes in bird-plant networks with 7 or more occurrences. Bottom row represent nodes in insect-plant networks with 11 or more occurrences. Left columns represent unstandardized local closeness centrality values. Right column represents the same points after z-score standardization with respect to 1000 null model randomization (see text).}
    \label{fig:S3}
  \end{center}
\end{figure}

\begin{figure}[h!]
\centering
    \includegraphics[width=\columnwidth]{figs/pdf/sup/MetawebCentralityPlot.pdf}%
    \caption{Species-level closeness centrality in the global metaweb as a function of  environmental niche breadth (A,C) or geographic range size (B,D). Plots in the left column show observed centrality values(A,B), while those in the right column show $z-$transformed centrality values based on permutations of local networks(C,D). Points and line colors respond to plant (green) or floral visitor (purple) nodes, respectively. Colored lines depict linear relationships between variables with $POP_2$ values $<0.05$ in our Bayesian GLMs (included for illustrative purposes only).}
    \label{fig:metawebScale}
\end{figure}



\begin{figure}[h!]
  \begin{center}
    \includegraphics[width=\textwidth]{figs/pdf/sup/Cooccurence_vs_GeographicRange.pdf}
    \caption{Relationship between log geographic range size (x-axis), and number of potential partners (extraguild species) it cooccurs with in our dataset. Shaded lines represent 95\% confidence intervals of a linear model. For insect-pollinated plants, we see a significant positive relationship between a species' geographic range size and the diversity of insect species it coocures with alongside its range (alpha, p). We see no significant effect of range size on potential interactor richness for avian-pollinated plants, or either pollinator group}
    \label{fig:S4}
  \end{center}
\end{figure}

\begin{figure}[h!]
  \begin{center}
    \includegraphics[width=\textwidth]{figs/pdf/sup/closenesscomparison.pdf}
    \caption{Comparison of unstandardized (x-axis) and z-standardized (y-axis) closeness centrality in global metawebs across our four node types. Points are colored by interaction turnover across sites. While there is a weak positive relationship between our centrality measure across most node types, the lack of strong explanatory power indicated that our z-transformation procedure does indeed capture alternative aspects of node connectivity at the metaweb scale.}
    \label{fig:S5}
  \end{center}
\end{figure}

\clearpage
Below we present diagnostic pairs plots of all bayesian multivariate regression models presented in the main text. Despite a positive correlation between log geographic range size and log environmental range area (Pearson's $r=0.56$) and negative correlation between fit slope values in some models, correlation structure did not prevent effective inference or MCMC chain convergence ($\hat{R}$ <1.01 across all presented models). 
\begin{figure}[h!]
  \begin{center}
    \includegraphics[width=\textwidth]{figs/pdf/sup/Pairplot_Insect_Visitors_Metaweb.pdf}
    \caption{Pairs plots of fit slope and intercept parameters for models of closeness centrality of \textbf{insect floral visitors} in the \textbf{global metaweb}. Hex plots represent the bivariate density of parameter estimates sampled from the posterior, while histograms represent univariate densities of each parameter. Top plots represent a model using raw closeness centrality values as a response variable; lower plots represent a model predicting z-standardized centrality. While there is some negative slope correlation between slope parameters for niche breadth and geographic range, observed $\hat{R}$ values <1.01 indicate adequate chain convergence for inference.}    \label{fig:S6}
  \end{center}
\end{figure}

\begin{figure}[h!]
  \begin{center}
    \includegraphics[width=\textwidth]{figs/pdf/sup/Pairplot_Insect_Plants_Metaweb.pdf}
    \caption{Pairs plots of fit slope and intercept parameters for models of closeness centrality of \textbf{insect-visited plants} in the \textbf{global metaweb}. Hex plots represent the bivariate density of parameter estimates sampled from the posterior, while histograms represent univariate densities of each parameter. Top plots represent a model using raw closeness centrality values as a response variable; lower plots represent a model predicting z-standardized centrality. While there is some negative slope correlation between slope parameters for niche breadth and geographic range, observed $\hat{R}$ values <1.01 indicate adequate chain convergence for inference.}
    \label{fig:S7}
  \end{center}
\end{figure}

\begin{figure}[h!]
  \begin{center}
    \includegraphics[width=\textwidth]{figs/pdf/sup/Pairplot_Avian_Visitors_Metaweb.pdf}
    \caption{Pairs plots of fit slope and intercept parameters for models of closeness centrality of \textbf{avian floral visitor} in the \textbf{global metaweb}. Hex plots represent the bivariate density of parameter estimates sampled from the posterior, while histograms represent univariate densities of each parameter. Top plots represent a model using raw closeness centrality values as a response variable; lower plots represent a model predicting z-standardized centrality.}
    \label{fig:S8}
  \end{center}
\end{figure}

\begin{figure}[h!]
  \begin{center}
    \includegraphics[width=\textwidth]{figs/pdf/sup/Pairplot_Avian_Plants_Metaweb.pdf}
    \caption{Pairs plots of fit slope and intercept parameters for models of closeness centrality of \textbf{avian-visited plants} in the \textbf{global metaweb}. Hex plots represent the bivariate density of parameter estimates sampled from the posterior, while histograms represent univariate densities of each parameter. Top plots represent a model using raw closeness centrality values as a response variable; lower plots represent a model predicting z-standardized centrality. While there is some negative slope correlation between slope parameters for niche breadth and geographic range, observed $\hat{R}$ values <1.01 indicate adequate chain convergence for inference.}
    \label{fig:S9}
  \end{center}
\end{figure}


\begin{figure}[h!]
  \begin{center}
    \includegraphics[width=\textwidth]{figs/pdf/sup/Pairplot_Plants_Local.pdf}
    \caption{Pairs plots of fit slope and intercept parameters for models of closeness centrality of \textbf{plants} in \textbf{local webs}. Hex plots represent the bivariate density of parameter estimates sampled from the posterior, while histograms represent univariate densities of each parameter. Top plots represent plants in avian-plant networks; lower plots represent plants in insect-plant networks. While there is some negative slope correlation between slope parameters for niche breadth and geographic range, observed $\hat{R}$ values <1.01 indicate adequate chain convergence for inference.}
    \label{fig:S10}
  \end{center}
\end{figure}

\begin{figure}[h!]
  \begin{center}
    \includegraphics[width=\textwidth]{figs/pdf/sup/Pairplot_Visitors_Local.pdf}
    \caption{Pairs plots of fit slope and intercept parameters for models of closeness centrality of \textbf{floral visitors} in \textbf{local webs}. Hex plots represent the bivariate density of parameter estimates sampled from the posterior, while histograms represent univariate densities of each parameter. Top plots represent birds in avian-plant networks; lower plots represent insects in insect-plant networks. While there is some negative slope correlation between slope parameters for niche breadth and geographic range in the model of local insect visitors, observed $\hat{R}$ values <1.01 indicate adequate chain convergence for inference.}
    \label{fig:S11}
  \end{center}
\end{figure}

\end{document}
